% Git setup commands for a LaTeX project
% This file documents the shell commands (not LaTeX code) for
% initializing a Git repository for a LaTeX book or thesis project.
% Run these commands in a terminal (PowerShell, bash, or cmd).

# ---- Step 1: Initialize the repository ----
# Navigate to your project folder, then:
git init

# ---- Step 2: Set your identity (run once globally) ----
git config --global user.name "Your Name"
git config --global user.email "you@example.com"

# ---- Step 3: Add a .gitignore file ----
# Copy the .gitignore from examples/ch09/gitignore-latex/ first!

# ---- Step 4: Add files and create the first commit ----
git add main.tex chapters/
git commit -m "Initial commit: book structure"

# ---- Step 5: Check status and history ----
git status
git log --oneline

# ---- Step 6: Connect to GitHub (optional) ----
# Create an empty repo on GitHub.com first, then:
git remote add origin https://github.com/username/repo-name.git
git branch -M main
git push -u origin main
